\documentclass[11pt]{article}
\usepackage[margin=1in]{geometry}
\usepackage[T1]{fontenc}
\usepackage{lmodern,microtype}
\usepackage{amsmath,amssymb,amsthm,hyperref}
\setlength{\emergencystretch}{2em}
\newtheorem{theorem}{Theorem}
\newtheorem{lemma}{Lemma}
\newcommand{\E}{\mathbb E}
\newcommand{\Prb}{\mathbb P}
\newcommand{\bits}{\{0,1\}}
\newcommand{\corr}{\operatorname{corr}}
\title{An explicit rate from the depth-three circuit construction}
\author{Tom Ayre}
\date{October 6, 2026}
\begin{document}
\maketitle

\noindent\textbf{Status.}
This note gives a mathematical proof of a quantitative strengthening of a
result in OpenAI's mathematics repository. Its clause-count-independent
restriction lemma has now been audited independently and is proved in full
below; the theorem is no longer conditional on that lemma. The external
input is the published quantitative, disjoint version of CNF sparsification.
The estimates needed for a growing proof parameter are checked directly.
This is not a claim of publication priority or external peer review.
Neither this strengthened theorem nor its new parameter schedule has been
checked in Lean. The accompanying paper provides the revised presentation.
This extended note was prepared with assistance from Codex.

\section{Statement and provenance}

Let $L$ be the language in Section 4 of
\emph{Beyond the Square-Root Exponent for Depth-Three Boolean Circuits},
September 23, 2026, and let $f_n$ be its $n$-bit membership function.
Let $S_3(f_n)$ denote the minimum total number of gates in an exact
unbounded-fan-in OR--AND--OR circuit for $f_n$, with free input negations,
constants, and arbitrary sharing, as in that manuscript.

The repository states $\log_2 S_3(f_n)/\sqrt n\to\infty$.
The quantitative sparsification input gives the following stronger rate.

\begin{theorem}[Quantitative strengthening]
There are absolute constants $c>0$ and $n_0$ such that
\[
 S_3(f_n)>2^{c\sqrt{n\ln n/\ln\ln n}}
 \qquad(n\ge n_0).
\]
The language and its polynomial-time membership algorithm are unchanged.
For any absolute sparsification constant $K$ satisfying (S) below, one may take
$c=1/(6144\sqrt K)$.
\end{theorem}

The source revision inspected was
\texttt{adc7f1241b42e322a6451854ab7e4b4c146bf78a}.
The relevant manuscript sources are in the
\href{https://github.com/openai/math/tree/adc7f1241b42e322a6451854ab7e4b4c146bf78a/preprints/Beyond-the-Square-Root-Exponent-for-Depth-Three-Boolean-Circuits-September-23-2026/build/latex/sections}{original manuscript's LaTeX section directory}
and are named
\path{02_restriction.tex}, \path{03_sparse_partition.tex},
\path{04_construction.tex}, and \path{05_lower_bound.tex}.
The existing Lean files \path{SparseTailBounds.lean},
\path{UniformSparseBounds.lean}, and \path{CircuitQuantitative.lean}
in \path{lean/OAI/Computability/DepthThree/}
also state finite numerical inequalities relevant to this derivation.
Their fixed-parameter eventual bounds are not being applied to parameters
that grow with $n$; the finite inequalities are rechecked below instead.

\section{Two inputs}

For a real function $G$ on a Boolean cube, define
\[
 \corr_b(G)=\max_{J}\left|\E[GJ]\right|,
\]
where $J$ ranges over indicators of all CNFs of width at most $b$,
with no clause-count restriction. Constants are allowed.

\paragraph{R: The repository restriction lemma, audited below.}
\begin{lemma}\label{lem:restriction}
Let $H$ be a width-$k$ CNF on $d$ variables. Independently retain each
variable with probability $0<p<1$, obtaining a live set $T$, and assign the
remaining variables uniformly, obtaining $\rho$. If
\[
 \theta=\sum_{\ell=b+1}^{k}\binom{k}{\ell}
                   \left(\frac{2p}{1-p}\right)^\ell<1,
\]
then there is a nonnegative random variable $W_H=W(H,T,\rho)$ such that
\[
 \E W_H\le(1-\theta)^{-1},\qquad
 |\E_z G(z)H_{T,\rho}(z)|\le W_H\corr_b(G)
\]
for every real function $G$ on the restricted cube. Here $k\ge0$ and
$b\ge1$ are integers. All parameters in this statement are finite; it has
no fixed-width asymptotic quantifier.
\end{lemma}

\paragraph{S: Quantitative disjoint sparsification.}
There is an absolute constant
$K\ge2$ such that every width-$b$ CNF indicator on $m\ge1$ variables is
a sum of at most $2^{m/8}$ pairwise disjoint width-$b$ CNF indicators,
each with at most $M(b)m$ clauses, where $M(b)$ is an integer and
\begin{equation}\label{eq:sparse}
 M(b)\le(Kb)^{Kb}\qquad(b\ge2).
\end{equation}
We use the quantitative counting form stated by
Williams~\cite[Section~4, p.~6:9]{Williams}, combining
Calabro, Impagliazzo, and Paturi~\cite[Section~6]{CIP} with
Dell et al.~\cite[Appendix~A, Lemma~A.1 and its following proof]{DHMTW}.
Fix an absolute error $\eta_0\in(0,1/8]$ in that theorem. Its quantitative
estimate means that the density before adding branch-defining units is
at most $D(b)\le A(b/\eta_0)^{Cb}$ for absolute $A,C\ge1$ independent
of $b$ and $m$. The finite recursion has at most $2^{\eta_0m}$ leaves
for every $m$; the finite quantifiers are explicit in Lemma~A.1
of~\cite{DHMTW}, rather than an asymptotic running-time convention.

For clarity, at a flower whose clauses are $H\vee P_i$, the heart branch
enforces $H$, while the petal branch enforces all $P_i$ and the unit clauses
negating each literal of $H$. These branches partition the solutions.
Along a consistent leaf at most one distinct added unit per variable is
needed, so at most $m$ clauses are added; discard inconsistent leaves.
The width is preserved. The reduction step removes clauses properly
containing a unit, so these units create no new flowers or leaves, as
explained in~\cite[Appendix~A]{DHMTW}. Induction gives a disjoint partition.
Set $M(b)=\lceil D(b)\rceil+1$ and choose any absolute
$K\ge\max\{2,2C,\eta_0^{-1},3A\}$. Then, for every $b\ge2$,
\[
 M(b)\le3A(b/\eta_0)^{Cb}\le(Kb)^{Kb},\qquad
 2^{\eta_0m}\le2^{m/8}.
\]
This proves (S) uniformly for every finite $b,m$, including $b>m$.
No numerical value of $K$ is claimed or required.

This replaces the repository's particular sparsification constant,
whose statement supplies existence for each fixed $b$ but not the
single-exponential-in-$b\log b$ estimate needed here.

\section{Proof of the restriction lemma}

This section reconstructs the source argument and checks the two essential
claims separately: the finite signed expansion, and its expected absolute
coefficient mass. In particular, no independence between a denominator and
the restriction is assumed.

\begin{proof}[Proof of Lemma~\ref{lem:restriction}]
Normalize $H$ by deleting tautological and true clauses and deleting repeated
literals within each remaining clause. Keep clause indices, including
repeated clauses. Write $E_i$ for the scope of clause $C_i$. A violation of
$C_i$ specifies a unique pattern on $E_i$; the false clause has empty scope.
Fix a restriction $(T,\rho)$, and let $v_i(z)$ be its violation indicator.
For $U\subseteq T$ set
\[
 q=\sum_i v_i,\qquad
 a_U=\sum_{i:\,|E_i\cap(T\setminus U)|\le b}v_i,\qquad
 c_U=\sum_{i:\,E_i\cap(T\setminus U)=\varnothing}v_i.
\]
Call an index hard at $U$ if $|E_i\cap(T\setminus U)|>b$.
The function $K_0=\mathbf1[a_\varnothing=0]$ is a width-$b$ CNF.

A path $P=(i_1,\ldots,i_j)$ starts at $U_0=\varnothing$, selects an index
hard at $U_{h-1}$, and sets $U_h=U_{h-1}\cup(E_{i_h}\cap T)$.
Each step adds at least $b+1$ coordinates, so paths have length at most
$L=\lfloor |T|/(b+1)\rfloor$. On the support of
$\prod_{h=1}^j v_{i_h}$, the selected indices are distinct and
$a_{U_h}\ge c_{U_h}\ge h$.
All fractions below are defined as zero when their numerator vanishes.

For a nonempty path define
\[
 r_P=\frac{\prod_{h=1}^jv_{i_h}}
                {q\prod_{h=1}^{j-1}a_{U_h}},\qquad
 t_P=\frac{\prod_{h=1}^jv_{i_h}}
                {\prod_{h=1}^{j}a_{U_h}}.
\]
Since $q=a_{U_j}+\sum_{i\text{ hard at }U_j}v_i$,
\[
 r_P=t_P-\sum_{i\text{ hard at }U_j}r_{Pi}.
\]
At a point with $K_0=1$ and $q>0$, all violated clauses are initially hard,
and $\sum_{|P|=1}r_P=1$. Iterating the displayed identity until there are
no further paths gives
\begin{equation}\label{eq:audited-expansion}
 H_{T,\rho}=K_0+
   \sum_{j=1}^{L}(-1)^j\sum_{|P|=j}K_0t_P.
\end{equation}
At points with $K_0=0$ every term vanishes; at points with $q=0$ the
nonempty-path terms vanish and $K_0=H_{T,\rho}=1$.
Thus this is an identity on the whole cube, not just on its error set.

Call a path compatible with $\rho$ when its selected violating patterns
are mutually consistent and agree with $\rho$ on fixed coordinates.
For such a path, all selected violations prescribe a single pattern
$\sigma_P$ on $U_j$. Denote its cylinder by $Z_P$.
The numerator of $t_P$ is $\mathbf1_{Z_P}$.
On $Z_P$ the function $c_{U_h}$ is a positive constant $c_h$, because all
variables of every counted clause have been fixed. Moreover,
\[
 a_{U_h}=c_h+w_h,
 \qquad
 w_h=\sum_{i:\,0<|E_i\cap(T\setminus U_h)|\le b}v_i.
\]
After substituting the final pattern $\sigma_P$, the indexed clauses
counted by $w_h$ have width at most $b$ on the remaining coordinates.
Their multiplicities and any constant clauses are retained.

For any indexed family of clauses and positive integer $c$, randomly
permute its clause objects together with $c$ distinct markers, and take
the conjunction of the clauses preceding the first marker. On an assignment
violating $w$ of those clauses, this conjunction holds with probability
$c/(c+w)$: the earliest of the $w$ violated clause objects and $c$ markers
must be a marker. This is a finite convex mixture of CNF indicators.
Apply the construction independently for each $h$, then conjoin the sampled
CNFs with $K_0$ and the unit clauses defining $Z_P$. Every resulting test
$J_P$ has width at most $b$, using $b\ge1$. It satisfies
\[
 K_0t_P=\omega(P)\,\E J_P,
 \qquad \omega(P)=\prod_{h=1}^j\frac1{c_h}.
\]
An incompatible path contributes zero. Define
\[
 W_0=1,\qquad
 W_j=\sum_{\substack{|P|=j\\P\text{ compatible with }\rho}}\omega(P),
 \qquad W=\sum_{j=0}^{\lfloor d/(b+1)\rfloor}W_j.
\]
Equation~\eqref{eq:audited-expansion} and the triangle inequality prove
$|\E_zGH_{T,\rho}|\le W\corr_b(G)$ for every $G$ simultaneously.
The construction of $W$ never uses $G$.

It remains to prove $\E W_j\le\theta^j$ for $j\ge1$.
Generate a uniform full assignment $y$, independently of $T$, and set
$\rho=y|_{[d]\setminus T}$. Conditional on $(T,\rho)$, a compatible path
has all its clauses violated by $y$ with probability $2^{-|U_j|}$.
Consequently
\begin{equation}\label{eq:audited-reweight}
 \E W_j=\E_{T,y}
 \sum_{\substack{|P|=j\\C_{i_1},\ldots,C_{i_j}\text{ violated by }y}}
 2^{|U_j|}\prod_{h=1}^j\frac1{c_{U_h}(y|_T)}.
\end{equation}
This is conditional cylinder averaging; it makes no claim of independence
between the path and its denominators.

Fix $y$, and for a remainder $R\subseteq[d]$ put
\[
 V_R=\{i:C_i(y)=0,\ E_i\cap R=\varnothing\}.
\]
An actual path has $R_h=T\setminus U_h$ and
$D_h=U_h\setminus U_{h-1}$, with
\[
 i_h\in V_{R_h},\qquad D_h\subseteq E_{i_h},\qquad |D_h|>b,
 \qquad R_{h-1}=R_h\mathbin{\dot\cup}D_h,\qquad
 c_{U_h}(y|_T)=|V_{R_h}|.
\]
Its terminal remainder $R_j$ and the reverse sequence
$(i_j,D_j),\ldots,(i_1,D_1)$ determine both $T$ and its entire indexed
path. Hence this is an injective encoding even when clauses repeat.

Let $\alpha=2p/(1-p)$ and
$\pi(R)=p^{|R|}(1-p)^{d-|R|}$. The weight of an encoded path in
\eqref{eq:audited-reweight}, including its live-set probability, is exactly
\[
 \pi(T)2^{|U_j|}\prod_h\frac1{|V_{R_h}|}
   =\pi(R_j)\prod_h\frac{\alpha^{|D_h|}}{|V_{R_h}|}.
\]
Enlarge the reverse sequences by allowing, at any state $R$, every choice
$i\in V_R$ and $D\subseteq E_i$ with $|D|>b$, followed by $R\cup D$.
Here $D\cap R=\varnothing$ because $i\in V_R$.
All weights are nonnegative, and the outgoing weight at each state with
$V_R\ne\varnothing$ is at most
\[
 \frac1{|V_R|}\sum_{i\in V_R}
          \sum_{\ell=b+1}^{|E_i|}\binom{|E_i|}{\ell}\alpha^\ell
 \le\sum_{\ell=b+1}^k\binom{k}{\ell}\alpha^\ell=\theta.
\]
If $V_R$ is empty there are no outgoing steps. Induction on the number
of reverse steps bounds their total weight from any state by $\theta^j$.
Summing over terminal remainders costs only
$\sum_R\pi(R)=1$. Thus the expected path mass over $T$, computed using
\eqref{eq:audited-reweight}, is at most $\theta^j$ for every fixed $y$, and
therefore $\E W_j\le\theta^j$.

Finally $\E W\le\sum_{j=0}^{\lfloor d/(b+1)\rfloor}\theta^j
\le(1-\theta)^{-1}$. Empty formulas, false clauses, repeated clauses,
empty live sets, and $k\le b$ all satisfy the same argument.
\end{proof}

\section{The fixed language and finite moment estimates}

For input length $n$, the construction uses
\[
 d=\lfloor n/5\rfloor,\quad r=\lceil d^{2/3}\rceil,\quad
 t=\max\{j\ge0:j\text{ even},\ j^6\le d\}.
\]
For $d\ge4096$,
\begin{equation}\label{eq:rt}
 r\ge d^{2/3},\qquad \tfrac12d^{1/6}\le t\le d^{1/6}.
\end{equation}
For $d\ge64$ we already have $t\ge2$, which suffices for the finite
moment estimate; the stronger lower bound in~\eqref{eq:rt} is only used
eventually.
Its input contains a $d$-bit data string $x$, a linear-hash seed $u$,
a monic degree-$r$ binary polynomial $P$, and $t$ elements
$\beta_0,\ldots,\beta_{t-1}$ of $\mathbb F_2[Z]/(P)$.
Membership tests whether
\[
 \lambda\!\left(\sum_{j=0}^{t-1}\beta_jh_u(x)^j\right)=0,
\]
where $\lambda$ selects the constant coefficient of the reduced
representative. The input blocks fit and the evaluator runs in polynomial
time as proved in the source. No parameter chosen below changes it.

Choose $P$ irreducible, the hash seed uniformly, and the coefficients
uniformly in $\mathbb F_{2^r}$. Write $g=+1$ on acceptance and $g=-1$
on rejection for the resulting data slice. The following are finite estimates from
Section 4 of the source, not its asymptotic corollaries.

On any fixed restricted data cube of dimension $m$:
\begin{itemize}
\item The hash is noninjective with probability at most $2^{m-r}$.
\item Conditional on an injective hash, the signs of the membership
function are unbiased and $t$-wise independent across distinct cube points.
\item Conditional on such a hash, for any one CNF indicator $J$, the probability that
$|\E_z g(z)J(z)|>2^{-m/4}$ is at most $t^t2^{-mt/4}$.
\end{itemize}
The third estimate follows by expanding the $t$th moment. Any term with
a point appearing exactly once has expectation zero. The remaining
ordered tuples number at most $t^t2^{mt/2}$. Markov's inequality gives
the stated tail.

There are at most
\[
 2^{M m[1+b\log_2(2m+1)]}
\]
CNF indicators with at most $Mm$ clauses of width at most $b$.
Thus the union-bound probability is
\[
 2^{Mm[1+b\log_2(2m+1)]+t\log_2t-mt/4}\le2^{-mt/8},
\]
provided
\begin{equation}\label{eq:union}
 \frac{M[1+b\log_2(2m+1)]}{t}
       +\frac{\log_2t}{m}\le\frac18.
\end{equation}
Disjoint sparsification then bounds correlation against every width-$b$
CNF by $2^{m/8}2^{-m/4}=2^{-m/8}$ simultaneously.

\section{Growing parameters and their verification}

All unqualified logarithms in the parameter choices are natural.
For sufficiently large $d$, set
\begin{align*}
 b&=\left\lfloor\frac{\ln d}{24K\ln\ln d}\right\rfloor,&
 s&=\frac{\sqrt b}{256},& B&=64(s+1),\\
 p&=\frac{B}{\sqrt d},& \mu&=B\sqrt d,&
 k&=\left\lceil3(s+1)\sqrt d\right\rceil.
\end{align*}
We work beyond a fixed absolute threshold so that $b\ge65536$, $s\ge1$,
and $p\le1/2$. Let $M=M(b)$ from (S).

\paragraph{The sparse-test union bound.}
Eventually $\ln(Kb)\le2\ln\ln d$, so \eqref{eq:sparse} gives
\[
 \ln M\le Kb\ln(Kb)\le\frac{\ln d}{12},\qquad M\le d^{1/12}.
\]
Uniformly for $\mu/2\le m\le2\mu$, we have
$\log_2(2m+1)=O(\ln d)$ and $b=O(\ln d)$.
More explicitly, \eqref{eq:rt} gives the finite envelope
\begin{align*}
 \frac{M[1+b\log_2(2m+1)]}{t}+\frac{\log_2t}{m}
 &\le 2M d^{-1/6}
       [1+b\log_2((4B+1)\sqrt d)]
       +\frac{\ln d}{3B\ln2\,\sqrt d}\\
 &\le 2d^{-1/12}
       [1+b\log_2((4B+1)\sqrt d)]
       +\frac{\ln d}{3B\ln2\,\sqrt d}.
\end{align*}
Consequently,
\[
 \frac{M[1+b\log_2(2m+1)]}{t}
       +\frac{\log_2t}{m}
 =O\!\left(d^{-1/12}(\ln d)^2
                   +\frac{\ln d}{\sqrt d}\right)=o(1).
\]
The constants here are absolute once the absolute $K$ is fixed.
This proves \eqref{eq:union} uniformly over the entire window.

\paragraph{Hash and restriction tails.}
Since $B=O(\sqrt{\ln d})$, we have $r/\mu\to\infty$.
In particular, eventually $r\ge(33/16)\mu$. Within the window,
\[
 2^{m-r}\le2^{-\mu/16},\qquad
 2^{-mt/8}\le2^{-\mu/16}.
\]
For the binomial random variable $m=|T|$ of mean $\mu$, the source's
elementary exponential-moment bounds give
\[
 \Prb[m\notin[\mu/2,2\mu]]\le2\cdot2^{-\mu/16}.
\]
Averaging the joint failure probability over the parameter bits therefore
gives one setting of those bits, for this length, whose sign function $g$
satisfies
\begin{equation}\label{eq:good}
 \Prb\left[m\notin[\mu/2,2\mu]
       \text{ or }\corr_b(g_{T,\rho})>2^{-m/8}\right]
       \le4\cdot2^{-\mu/16}.
\end{equation}
This is a pointwise existence argument at each $d$; it does not exchange
the original fixed-parameter asymptotic quantifiers.

\paragraph{The restriction lemma still applies.}
For $d\ge1$,
\[
 k\le[3(s+1)+1]\sqrt d,\qquad
 \frac{2pk}{1-p}\le4B[3(s+1)+1]=:Q.
\]
Using $B=64(s+1)$ and $s\ge1$,
\[
 Q\le1024(s+1)^2\le4096s^2=\frac b{16}.
\]
By $\binom{k}{\ell}\le k^\ell/\ell!$ and
$\ell!\ge(\ell/e)^\ell$,
\[
 \theta\le\sum_{\ell>b}\frac{Q^\ell}{\ell!}
 \le\sum_{\ell>b}(e/16)^\ell<\frac12.
\]
Thus (R) supplies $\E W_H\le2$ for every width-$k$ test $H$.

\section{The circuit contradiction}

On the good restrictions in \eqref{eq:good}, the restricted correlation
is at most $W_H2^{-\mu/16}$. On the bad restrictions it is at most one.
Uniform completion of a random restriction is uniform on $\bits^d$.
Taking expectations and using the triangle inequality gives
\begin{equation}\label{eq:corr}
 |\E_x g(x)H(x)|\le6\cdot2^{-\mu/16}
\end{equation}
for every width-$k$ CNF $H$. Independence of $W_H$ and the good event
is not required. Applying \eqref{eq:corr} to the constant-true CNF
shows that $\Prb[g=1]\ge1/4$ for sufficiently large $d$.

Suppose there were a circuit for $f_n$ with at most $S\le2^{s\sqrt d}$
total gates. Fix its parameter inputs to the setting chosen above.
Each middle AND gate is a CNF $H$ with at most $S$ clauses, and there
are at most $S$ such middle gates.
Delete from $H$ the clauses of width greater than $k$, obtaining $H'$.
Then $H'\ge H$ and
\[
 \E(H'-H)\le S2^{-k}.
\]
Every input accepted by $H$ is accepted by the top OR, so $gH=H$.
Equation \eqref{eq:corr} therefore implies
\[
 \E H\le6\cdot2^{-\mu/16}+S2^{-k}.
\]
Taking the union over middle gates yields
\begin{align*}
 \Prb[g=1]
 &\le6S2^{-\mu/16}+S^2 2^{-k}\\
 &\le6\cdot2^{-(3s+4)\sqrt d}
             +2^{-(s+3)\sqrt d}=o(1),
\end{align*}
contradicting $\Prb[g=1]\ge1/4$.
Hence $S_3(f_n)>2^{s\sqrt d}$ for every sufficiently large $n$.
Finally,
\[
s\sqrt d=\frac{\sqrt{bd}}{256}
       =\Theta\!\left(\sqrt{\frac{n\ln n}{\ln\ln n}}\right),
\]
because $d=\lfloor n/5\rfloor$ and
$b=\Theta(\ln d/\ln\ln d)$. More explicitly, eventually
\[
 d\ge\frac n6,\qquad \ln d\ge\frac{\ln n}{2},\qquad
 b\ge\frac{\ln d}{48K\ln\ln d}
   \ge\frac{\ln n}{96K\ln\ln n}.
\]
It follows that
\[
 s\sqrt d\ge\frac1{6144\sqrt K}
                 \sqrt{\frac{n\ln n}{\ln\ln n}},
\]
which proves the theorem with the announced absolute constant.

\section{What has and has not been established}

The contribution is an explicit rate, obtained by replacing a
fixed-width sparsification input with a quantitative one and checking
all subsequent estimates for a growing width. It is stronger than the
source's stated $2^{\omega(\sqrt n)}$ guarantee, which gives no specific
diverging factor. It does not reach $2^{n^{1/2+\epsilon}}$ for any fixed
$\epsilon>0$.

The displayed schedule has an enormous threshold. The condition
$b\ge65536$ alone requires
$\ln d/\ln\ln d\ge24K\cdot65536\ge3{,}145{,}728$; on the large-$d$
branch, $d$ has more than 24 million decimal digits. Other eventual
conditions may increase the threshold. These conservative choices say
nothing about the smallest lengths at which the lower bound is true.

The available density estimate $M(b)\le\exp(O(b\ln b))$ certifies a
polynomial density in $d$ when $b\ln b=O(\ln d)$, accounting for the
$\ln\ln n$ loss. A uniform disjoint bound $M(b)\le2^{Cb}$ with an
absolute $C$ and at most $2^{m/8}$ pieces would let us take
$b=\delta\ln d$ for a sufficiently small constant $\delta>0$, giving
$2^{\Omega(\sqrt{n\ln n})}$ by the same proof. This would be a result
for this language, not for majority. Section~7 of~\cite{CIP} asks about
a corresponding density improvement; no impossibility of the needed
disjoint bound is asserted here. Raising $t$ to another fixed power
of $d$, if the input budget and hash conditions still hold, would change
constants but not the rate obtained from the current density estimate.
Since $t$ specifies the coefficient blocks, that would also change the
language itself.

The independent audit of (R) found no gap; its complete finite proof is now
included above. In particular, the clause-count cancellation is justified
by the reverse encoding and its exact probability reweighting, rather than
an independence assumption. Separate reviews checked the quantitative
parameter calculation and the existing Lean restriction source.
The theorem is thus proved here in ordinary mathematics using the published
sparsification theorem, without an additional conjectural restriction
hypothesis. Novelty and external peer review remain separate questions.

The original
\href{https://github.com/openai/math/blob/adc7f1241b42e322a6451854ab7e4b4c146bf78a/lean/OAI/Computability/DepthThree/Restriction.lean}{Lean restriction file}
contains theorem \path{cnf_restriction_bound}, whose hypotheses are those
of (R); it does not assume a certificate for the core bound. A source scan
of its local import closure found no \texttt{sorry}, \texttt{admit}, or
new \texttt{axiom} declarations. This is a source audit, not a fresh kernel
check: no Lean or Lake executable, dependency directory, or compiled build
artifacts were available locally. Formal integration of the quantitative
disjoint sparsification input and the growing schedule remains to be done.

The accompanying \texttt{check\_quantitative\_depth\_three.py} checks the
elementary polynomial identities and factorial-tail constant with exact
arithmetic, and checks nine illustrative parameter choices in the log
domain. The sample values of $K$ are not asserted to be certified
sparsification constants. The separate
\texttt{check\_restriction\_lemma.py} checks the finite restriction
identities with exact rational arithmetic on small formulas, including
overlaps, repeated clauses, and constant cases. Its completed run covered
all 512 clause subsets on two variables and 11 additional cases at two
rational probabilities: 534 formula--parameter cases, 11,254 restrictions,
and 43,524 completions. All checks passed, including nonzero three-step
paths and finite mixtures with seven indexed clauses. These diagnostics supplement
the mathematical proofs; they do not prove their universal statements.

\section{Novelty audit: October 6, 2026}

\paragraph{Conclusion and scope.}
A targeted primary-source literature audit found no earlier statement of
this explicit rate for the repository language, and no independent result
that subsumes the claimed bound for a polynomial-time language in the same
unrestricted OR--AND--OR gate model. The supported assessment is
\emph{an apparently unstated quantitative corollary of the repository
construction}. This is evidence for novelty of the rate, not a certificate
of publication priority, an assessment of journal significance, or proof
that no unpublished derivation exists.

The comparison requires the same resource: total AND/OR gates, arbitrary
bottom fan-in and sharing, free input negations, and the total number of
input bits as the size parameter. Larger bounds for restricted circuits
do not establish the unrestricted statement. Neither do existence bounds
for nonexplicit truth tables establish a bound for a fixed language in P.

\paragraph{What is already in the repository.}
The source introduction states only
$\log_2 S_3(f_n)/\sqrt n\to\infty$. Its lower-bound proof expressly fixes
$A$, then $B,b$, before letting $n$ grow. Its sparsification lemma gives
an existence statement for $M(b,\eta)$ and a recursive construction,
without the estimate $(Kb)^{Kb}$ used here. Searches of the relevant
manuscript sections, bibliography, DepthThree Lean modules, final theorem,
and formal scope note found neither the present rate nor this growing-width
schedule. The public scope note is
\url{https://raw.githubusercontent.com/openai/math/main/lean/docs/112.md}.

The language, restriction lemma, signed expansion, marker mixtures, and
backwards clause-count cancellation are all ingredients of the source.
The independent audit above validates those ingredients; it does not claim
to have invented them. Quantitative disjoint sparsification is also prior
work \cite{CIP,DHMTW,Williams}. The contribution here is its combination
with a growing parameter choice and the finite estimates needed to justify
the resulting explicit rate.

\paragraph{Closest external comparisons.}
\begin{itemize}
\item G\"o\"os, Guan, and Mosnoi, MFCS 2023, describe exceeding the
square-root exponent for unrestricted explicit depth-three functions as
open. Their main exponential inner-product bound concerns bottom fan-in
two \cite[pp.~51:1--51:2]{GGM23}.
\item Gurumukhani et al., February 3, 2026, again describe the unrestricted
$2^{\omega(\sqrt n)}$ bound as open. Their proved majority result is for
monotone circuits with bottom fan-in three. The stronger unrestricted
majority rate $2^{\Omega(\sqrt{n\log n})}$ appears as a consequence of
a conjecture, not as an achieved theorem \cite[Section~1]{GKKPRST}.
The 2024 local-enumeration paper similarly gives a conditional route and
an achieved bottom-fan-in-three result \cite{GPPST24}.
\item The February 12, 2026 version of \emph{Optimal Depth-Three Circuits
for Inner Product} matches an exponential lower bound by an upper bound
in the bottom-fan-in-two model; it does not supply an unrestricted
improvement \cite{GKPRT26IP}.
\item Goldreich and Tal's $\exp(\Omega(n^{3/5}))$ bound applies to
\emph{canonical} depth-three circuits, a restricted class, and therefore
does not subsume the present theorem \cite[Corollary~1.4]{GT15}.
\item Lin's September 4, 2026 revision gives unrestricted top-disjunction
depth-three bounds for parameterized $k$-OV, $k$-XOR, and $k$-SUM
\cite{Lin26}. For XOR and SUM, its displayed bound is
$\exp(\Omega(k\log(n/k)))$ at bit width
$m=\Theta(k\log(en/k))$ per item. The total bit length is
$N=\Theta(nk\log(en/k))$. In the advertised growing regime
$k=n^{o(1)}$, the exponent scale $k\log(n/k)=N^{o(1)}$ is below
$\sqrt N$. This bit-length comparison is our inference from its
parameters, not a statement that $n$ itself counts input bits. The
theorem's additional polynomial host-size constraints also prevent
choosing $k$ comparable to $n$ to obtain the claimed rate here.
\end{itemize}

These comparisons support the absence of a matching result in the sources
examined. Statements in early-2026 papers that a problem was open establish
the authors' reported state of knowledge then; they do not establish that
no later or unindexed work exists.

\paragraph{Search coverage and priority limits.}
The search used the exact manuscript title, combinations of
``depth-three'', ``depth-3'', ``Boolean'', ``quantitative'', and
``sparsification'', variants of $\sqrt{n\log n/\log\log n}$, and targeted
arXiv and ECCC searches. It followed relevant primary papers and compared
their hypotheses, rather than treating search snippets as theorem statements.
Repository mirrors were identified as copies of the same source and were
not counted as independent confirmations.

The manuscript is dated September 23, 2026, but the inspected local history
contains only initial commit
\texttt{adc7f1241b42e322a6451854ab7e4b4c146bf78a}, timestamped
October 6, 2026, 14:58:50 at UTC$-07{:}00$. A manuscript date or Git
timestamp alone does not establish first public disclosure. The repository
is very recent, so search indexing is especially incomplete. No authors
were contacted, and no exhaustive citation database or unpublished-work
search was performed.

A defensible attribution is:
\begin{quote}
We derive an explicit $2^{\Omega(\sqrt{n\log n/\log\log n})}$ lower
bound for the language of OpenAI (2026), by combining its construction
and restriction lemma with published quantitative disjoint sparsification.
This rate is not stated in the source manuscript or its accompanying
formalization, and we did not find it in the literature examined.
\end{quote}

\begin{thebibliography}{99}
\bibitem{CIP}
Chris Calabro, Russell Impagliazzo, and Ramamohan Paturi.
\emph{A Duality between Clause Width and Clause Density for SAT}.
Section 6, Improvement of sparsification.
\url{https://cseweb.ucsd.edu/~ccalabro/cip.pdf}.

\bibitem{DHMTW}
Holger Dell, Thore Husfeldt, Daniel Marx, Nina Taslaman, and Martin Wahlen.
\emph{Exponential Time Complexity of the Permanent and the Tutte Polynomial}.
Full version, Appendix A, 2012; journal version, 2014.
\url{https://arxiv.org/html/1206.1775v1}.

\bibitem{Williams}
R. Ryan Williams.
\emph{Counting Solutions to Polynomial Systems via Reductions}.
SOSA 2018, Section 4, page 6:9.
\url{https://people.csail.mit.edu/rrw/counting-sosa.pdf}.

\bibitem{GKKPRST}
Mohit Gurumukhani, Daniel Kleber, Ramamohan Paturi, Christopher Rosin,
Michael Saks, and Navid Talebanfard.
\emph{Optimal Monotone Depth-Three Circuit Lower Bounds for Majority}.
arXiv:2601.04072v2, February 3, 2026. Introduction.
\url{https://arxiv.org/html/2601.04072v2}.

\bibitem{GGM23}
Mika G\"o\"os, Ziyi Guan, and Tiberiu Mosnoi.
\emph{Depth-3 Circuits for Inner Product}.
MFCS 2023, LIPIcs 272, article 51; journal version, Information and
Computation, 2024.
\url{https://drops.dagstuhl.de/opus/volltexte/2023/18585/pdf/LIPIcs-MFCS-2023-51.pdf}.

\bibitem{GPPST24}
Mohit Gurumukhani, Ramamohan Paturi, Pavel Pudlak, Michael Saks,
and Navid Talebanfard.
\emph{Local Enumeration and Majority Lower Bounds}.
arXiv:2403.09134v3, May 23, 2024; CCC 2024.
\url{https://arxiv.org/abs/2403.09134v3}.

\bibitem{GKPRT26IP}
Mohit Gurumukhani, Daniel Kleber, Ramamohan Paturi, Christopher Rosin,
and Navid Talebanfard.
\emph{Optimal Depth-Three Circuits for Inner Product}.
arXiv:2601.04446v2, February 12, 2026.
\url{https://arxiv.org/abs/2601.04446v2}.

\bibitem{GT15}
Oded Goldreich and Avishay Tal.
\emph{Matrix Rigidity of Random Toeplitz Matrices}.
ECCC TR15-079, 2015, Section 1.2 and Corollary 1.4.
\url{https://eccc.weizmann.ac.il/report/2015/079/download/}.

\bibitem{Lin26}
Haoxing Lin.
\emph{Fine-Grained $\mathrm{AC}^0$ Lower Bounds for $k$-OV, $k$-XOR,
and $k$-SUM via Colored Subgraph Isomorphism}.
ECCC TR26-139, revision 2, September 4, 2026, Theorem 6.8.
\url{https://eccc.weizmann.ac.il/report/2026/139/revision/2/download/}.
\end{thebibliography}
\end{document}
